\section{Actually acquired occupation and causal resolution}\label{sec:geometry}
The task criterion is useful only if the relevant geometry is itself acquired by the experiment. Fix a declared exploration and a common predictive realization when one is available. At designated scored calls let $p_t=\E[Y_t\mid H_t]$ be the conditional mean of a bounded finite-dimensional target, with $p_t\in K$ for a compact $K\subset\R^q$. In a robust family this formula is interpreted modelwise, and a common executable representation must be verified separately. Define the excursion occupation measures
\[
 \nu_i(A)=\E_i\sum_{t\text{ scored in excursion}}\one_{\{p_t\in A\}}.
\]
Preparation or other zero-score calls are omitted from this numerator but included in the duration $d_i$. The actual physical-time measure from initial label $i$ is
\begin{equation}\label{eq:occupation}
 \nu_\infty^i=\sum_C w_C(i)\frac{\sum_{j\in C}\pi_C(j)\nu_j}{\pi_Cd}.
\end{equation}
All quantization identities below concern a specified prepared law. In a robust family, modelwise quantization gives a lower bound, whereas an implementable upper must use one common codebook and recursion. The measure in~\eqref{eq:occupation} is a subprobability measure. Transient boundary occupation has zero asymptotic mass. Let $\nu_N^i$ be the expected empirical scored occupation divided by $N$ physical calls, and define $\nu_\beta^i$ with normalized-discount weights. These measures are neither preparation measures nor freely chosen volume measures.

\begin{theorem}[Task-visible acquired-law criterion]\label{thm:occupation}
Assume continuous finite responses and~\eqref{eq:UI}. Let the scored feature maps be such that finite responses of every bounded Lipschitz function of $p_t$ are continuous in $x$. Then the following are equivalent: continuity of $x\mapsto\nu_\infty^i$ in the bounded-Lipschitz metric for all initial labels; uniform convergence $\nu_N^i\to\nu_\infty^i$ in that metric; uniform discounted convergence; and Theorem~\ref{thm:transfer}'s task criterion on the compact family of normalized nonnegative Lipschitz tests of $K$. Its quantitative bounds use the supremum of the corresponding approximate-corrector prices.
\end{theorem}
\begin{proof}
Use $\mathcal F=\{f:K\to[0,1]:\Lip(f)\le1\}$, which is compact in the uniform norm by Arzel\`a--Ascoli. Treat $f$ as an additional parameter and score $f(p_t)$ only at designated calls. Formula~\eqref{eq:class-gain} is exactly~\eqref{eq:occupation} tested against $f$. Finite-response continuity is joint in $(x,f)$ because uniform changes in $f$ change normalized expected scores by at most their uniform norm. Continuity of the measure implies joint continuity of the gain; conversely continuity for all $f$ on compact $X\times\mathcal F$ gives bounded-Lipschitz continuity. Apply Theorem~\ref{thm:transfer} on this product. This argument retains the true mass of every stratum and does not require a smooth coordinate chart.
\end{proof}

For a finite positive measure $\nu$ on $K$, define
\begin{equation}\label{eq:quantization}
 q_J(\nu)=\inf_{c_1,\ldots,c_J\in\conv K}
       \int\min_{1\le j\le J}\norm{p-c_j}^2\,d\nu(p).
\end{equation}
In finite-dimensional compact support the infimum is attained. The equality-of-infima formulation remains meaningful in a general Hilbert space, where attainment requires an additional argument. Put $D_K=\operatorname{diam}(K)$.

\begin{lemma}[Projection, actual mass and quantization]\label{lem:quantization}
For a fixed exploration, the best checkpoint risk with $J$ fixed response values is its conditional-variance baseline plus $q_J$ of the actual scored occupation. Every online observer with at most $J$ such values has at least this risk, even when those labels are also used for control. Further,
\begin{equation}\label{eq:q-BL}
 |q_J(\nu)-q_J(\nu')|
 \le(D_K^2+2D_K)\,d_{\rm BL}(\nu,\nu'),
\end{equation}
with a metric using bounded-plus-Lipschitz norm at most one. If $\lambda\le\nu$ has mass $w>0$ and
\[
 \lambda(B(z,r))\le C r^d\qquad(z\in\conv K,\ 0<r\le r_0),
\]
then, whenever $r=(w/(2CJ))^{1/d}\le r_0$,
\begin{equation}\label{eq:smallball}
 q_J(\nu)\ge\frac w2\left(\frac{w}{2CJ}\right)^{2/d}.
\end{equation}
A global cover of $K$ by $J$ radius-$h_J$ balls gives $q_J(\nu)\le\nu(K)h_J^2$.
\end{lemma}
\begin{proof}
Conditional squared projection gives
$\E\norm{Y-c}^2=\E\norm{Y-p}^2+\E\norm{p-c}^2$ for each measurable response, with independent randomized choices included by conditioning. For a fixed codebook the pointwise nearest response is optimal. A checkpoint can make this assignment measurably; any online response is bounded below by the same pointwise distance. The function $\min_j\norm{p-c_j}^2$ has bound $D_K^2$ and Lipschitz constant at most $2D_K$. Taking infima proves~\eqref{eq:q-BL}. The union of $J$ radius-$r$ balls has $\lambda$-mass at most $JCr^d=w/2$; the remaining mass contributes at least $wr^2/2$. The cover bound is immediate. No mass is normalized away.
\end{proof}
The baseline may vary with the exploration. The pointwise lower remains valid when the observer jointly controls and predicts, but it is then the occupation law generated by \emph{that} observer that enters.

\subsection{A counted causal implementation}
A geometric cover alone is not an executable predictor. The following assumptions specify the missing causal content. Fix an exploration using $Q$ retained control labels. The predictive realization may use a state carrier larger than $K$: write its state as $s$, and require that the conditional task mean is $p=f(s)\in K$. Suppose it has a common known initialization at physical preparation, a jointly Borel update $F(s,a,y)$ and a task map $f$ with Lipschitz constant $L_f$. When the carrier is the conditional-mean space itself, take $f$ to be the identity. All quantization measures below are measures of $p=f(s)$, not of an unrelated carrier coordinate. A budget-independent relation between exact and candidate states is specified. For every $J$ there is a globally valid $J$-center projection of error at most $h_J$ on the full candidate domain; the update, projection and numerical perturbations preserve that relation. The report kernel is continuous on that domain in the declared task or total-variation modulus. Continuity is not a substitute for the following same-report stability estimate.

Along the \emph{actual} reports and exploration actions, suppose every admissible candidate recursion satisfies
\begin{equation}\label{eq:gain-recursion}
 e_t\le L_t e_{t-1}+h_J+\eta_{\rm num},
 \qquad
 G_t=1+L_tG_{t-1},\quad G_{\rm prep}=1,
 \qquad e_{\rm prep}\le h_J+\eta_{\rm num}.
\end{equation}
The predictor, but not the retained control label, is initialized at preparation. Assume
\begin{equation}\label{eq:occupied-stability}
 \sup_{x,N}\frac1N\E_x\sum_{t\le N,\ t\text{ scored}}G_t^2\le S_0<\infty.
\end{equation}
A limiting or horizon-dependent version gives the corresponding limiting or horizon-dependent bound. This is an actual-acquisition condition, not stability on a formal reachable set.

\begin{theorem}[Acquired geometry to causal memory]\label{thm:geometry}
Under these assumptions, a $QJ$-label observer implements the exploration and a predictor whose squared excess risk at every physical horizon is at most
\begin{equation}\label{eq:geometric-upper}
 \left[L_f(h_J+\eta_{\rm num})\sqrt{S_0}+\xi\right]^2,
\end{equation}
where $\xi$ bounds the numerical task readout error. Its retained control coordinate survives preparation. The checkpoint excess is exactly $q_J(\nu_N)$; every online $M$-label observer has excess at least $q_M(\nu_N)$ for its own actual law. Under Theorem~\ref{thm:occupation}, finite-acquisition and average quantization laws differ by~\eqref{eq:q-BL} and the task-corrector modulus.

For optimized average risk, suppose every legal $M$-label observer has baseline at least $B_*$ and its own actual occupation contains the same uniform small-ball witness parameters $(w,C,d,r_0)$. Suppose a fixed $Q$-label exploration attains $B_*$ and satisfies the global cover and stability assumptions with $h_J\le C_1J^{-1/d}$. Then, for $M\ge2Q$ and exact arithmetic,
\begin{equation}\label{eq:optimized-geometry}
 cM^{-2/d}\le\inf_{\gamma\in\Gamma_M}\calR_\infty(\gamma)-B_*
 \le C_2M^{-2/d}.
\end{equation}
The lower quantifiers are over all legal observers; the upper uses the displayed exploration.
\end{theorem}
\begin{proof}
Store $(q,c)$, the exploration label and the predictive center. Actions use only $q$, so the physical exploration law is unchanged. On an actual report, update the candidate from $c$ and project it; the next control label is updated by its original rule. There are exactly $QJ$ available pairs, including any internal clock already charged to the exploration. At preparation reset only the predictor coordinate to its common approximate prior. Induction in~\eqref{eq:gain-recursion} gives $e_t\le(h_J+\eta_{\rm num})G_t$. The task map and numerical readout yield pointwise root error at most $L_f(h_J+\eta_{\rm num})G_t+\xi$. Minkowski's inequality in the actual normalized scored occupation and~\eqref{eq:occupied-stability} give~\eqref{eq:geometric-upper}; its square, not its root, is the excess risk. Lemma~\ref{lem:quantization} and Theorem~\ref{thm:occupation} prove the checkpoint and temporal assertions.

For optimized matching, apply~\eqref{eq:smallball} to the law of each competing observer, using its at most $M$ task readouts even if the labels also implement control. This gives the uniform lower above $B_*$. For the upper choose $J=\lfloor M/Q\rfloor$, which is at least $M/(2Q)$, and use the attaining exploration with the counted product implementation. These are the extra assumptions required to pass from fixed-exploration geometry to optimized control; no such passage is automatic.
\end{proof}

\begin{proposition}[Survival-weighted expansion]\label{prop:survival}
Suppose an excursion survives each further scored step with conditional probability at most $s<1$, and every candidate-valid one-step Lipschitz multiplier is at most $L$. If $\sqrt{s}L<1$, then an actual-law stability constant as in~\eqref{eq:occupied-stability} is finite, bounded by a constant depending only on $s,L$ and the fixed number of non-scored preparation calls. Individual updates may expand, $L>1$.
\end{proposition}
\begin{proof}
At depth $k$, $G_k\le\sum_{j=0}^kL^j$. Minkowski in the weighted sequence space gives
\[
 \left(\sum_{k\ge0}s^kG_k^2\right)^{1/2}
 \le\sum_{j\ge0}L^j\left(\sum_{k\ge j}s^k\right)^{1/2}
 =\frac1{\sqrt{1-s}(1-\sqrt{s}L)}.
\]
Conditioning at each actual preparation and using at most one preparation per physical call bounds normalized cumulative occupation by this excursion sum, up to the fixed convention at preparation. No conditioning on successful survival removes its probability. Expected logarithmic contraction or random-product hypotheses may be weaker, but are not asserted here.
\end{proof}
This sufficient stability mechanism and the task-visible occupation criterion address different obstructions. The former realizes a recursive predictor; the latter permits temporal rank degeneration invisible to its test family. Neither implies the other without the displayed hypotheses.
